% =====================================================================
%  Q-scan morphology of scattered-light transients in LIGO O3 and O4a
%  Hardik Sher Singh, Thapar Institute of Engineering and Technology
%
%  Compile on Overleaf with pdfLaTeX:
%      pdflatex main -> bibtex main -> pdflatex main -> pdflatex main
%  (Overleaf's default recipe does this automatically.)
%
%  Figures are read from ../results/figures if the LaTeX source sits inside
%  the analysis repository, and from ./figures if the paper/ directory has
%  been uploaded to Overleaf on its own.  Both paths are declared below, so
%  neither arrangement needs an edit.
% =====================================================================

\documentclass[11pt,a4paper]{article}

% ---- layout ---------------------------------------------------------
\usepackage[margin=2.4cm]{geometry}
\usepackage{setspace}
\onehalfspacing
\usepackage{parskip}

% ---- encoding and fonts ---------------------------------------------
\usepackage[T1]{fontenc}
\usepackage[utf8]{inputenc}
\usepackage{lmodern}
\usepackage{microtype}

% ---- maths ----------------------------------------------------------
\usepackage{amsmath}
\usepackage{amssymb}
\usepackage{bm}
\usepackage{siunitx}
\sisetup{
  range-units      = single,
  separate-uncertainty = true,
  detect-weight    = true,
}

% ---- figures and tables ---------------------------------------------
\usepackage{graphicx}
\graphicspath{{../results/figures/}{figures/}{./}}
\usepackage{booktabs}
\usepackage{longtable}
\usepackage{array}
\usepackage[font=small,labelfont=bf]{caption}
\usepackage{subcaption}
\usepackage{float}

% ---- code -----------------------------------------------------------
\usepackage{xcolor}
\definecolor{codebg}{RGB}{247,247,247}
\definecolor{coderule}{RGB}{220,220,220}
\definecolor{keyword}{RGB}{0,90,150}
\definecolor{comment}{RGB}{110,110,110}
\usepackage{listings}
\lstset{
  basicstyle    = \ttfamily\footnotesize,
  keywordstyle  = \color{keyword}\bfseries,
  commentstyle  = \color{comment}\itshape,
  stringstyle   = \color{comment},
  backgroundcolor = \color{codebg},
  frame         = single,
  rulecolor     = \color{coderule},
  breaklines    = true,
  showstringspaces = false,
  columns       = flexible,
  xleftmargin   = 0.5em,
  language      = Python,
}

% ---- references and links -------------------------------------------
\usepackage[numbers,sort&compress]{natbib}
\bibliographystyle{unsrtnat}
\usepackage[hidelinks,breaklinks]{hyperref}
\hypersetup{
  pdftitle  = {Q-scan morphology of scattered-light transients in LIGO O3 and O4a},
  pdfauthor = {Hardik Sher Singh},
  colorlinks = true,
  linkcolor  = black,
  citecolor  = black,
  urlcolor   = blue!55!black,
}
\usepackage[nameinlink,capitalise]{cleveref}

% ---- draft flagging -------------------------------------------------
%  The results section carries a status notice, but that sits twenty pages in.
%  A reader who opens the PDF, reads the title and abstract, and skims would
%  see nothing at all -- which is exactly how a draft gets circulated as if it
%  were finished.  So the same condition that raises the status notice also
%  puts a banner on the title page, a marker in the PDF metadata, and a line
%  in the footer of every single page.  All four disappear together the moment
%  stage 06 writes numbers.tex; none of them can be forgotten separately.
\newif\ifresultspending
\IfFileExists{numbers.tex}{\resultspendingfalse}{\resultspendingtrue}

\ifresultspending
  \usepackage{fancyhdr}
  \pagestyle{fancy}
  \fancyhf{}
  \renewcommand{\headrulewidth}{0pt}
  \renewcommand{\footrulewidth}{0.4pt}
  \fancyfoot[L]{\footnotesize\textcolor{red!70!black}{\textbf{DRAFT}}}
  \fancyfoot[C]{\thepage}
  \fancyfoot[R]{\footnotesize\textcolor{red!70!black}{%
    results not yet computed}}
  \fancypagestyle{plain}{%
    \fancyhf{}%
    \renewcommand{\headrulewidth}{0pt}%
    \renewcommand{\footrulewidth}{0.4pt}%
    \fancyfoot[L]{\footnotesize\textcolor{red!70!black}{\textbf{DRAFT}}}%
    \fancyfoot[C]{\thepage}%
    \fancyfoot[R]{\footnotesize\textcolor{red!70!black}{%
      results not yet computed}}%
  }
  \hypersetup{
    pdftitle  = {[DRAFT -- results not yet computed]
                 Q-scan morphology of scattered-light transients in LIGO O3
                 and O4a},
    pdfsubject = {DRAFT. The analysis pipeline is validated but the production
                  run over the full O3 and O4a samples has not been executed;
                  the results section contains placeholders, not
                  measurements.},
  }
\fi

% ---- placeholders for numbers the pipeline has not yet produced -----
%  \result{...} renders in a visible colour so that no unfilled number can
%  reach a reader mistaken for a measurement.  Stage 06 of the pipeline
%  writes numbers.tex, which redefines these to the measured values.
\providecommand{\result}[1]{\textcolor{red!75!black}{\textbf{[#1]}}}

%  Every inline quantity quoted in the text is declared here with a visible
%  default.  Stage 06 writes numbers.tex, which \renewcommands them to the
%  measured values.  Declaring them here rather than only in numbers.tex means
%  a key the pipeline forgot to emit shows up as a red placeholder, and a key
%  it emits under the wrong name is a LaTeX error -- neither can pass silently.
\newcommand{\ResultStrongestFeature}{\result{feature}}
\newcommand{\ResultStrongestDelta}{\result{delta}}
\newcommand{\ResultStrongestChannel}{\result{channel}}
\newcommand{\ResultStrongestRho}{\result{rho}}
\newcommand{\ResultQuartileRatio}{\result{ratio}}
\newcommand{\ResultLagPeak}{\result{lag}}
\newcommand{\ResultClassifierAuc}{\result{AUC}}
\newcommand{\ResultClassifierEfficiency}{\result{efficiency}}
%  Names are exactly what gwglitch.latex.macro_name() produces from the
%  pipeline's keys: digits are spelled out because TeX forbids them in a
%  control sequence, so "o4a" becomes "OFoura".  Ugly, but mechanical and
%  therefore reliable; tests/test_latex.py checks the two stay in step.
\newcommand{\ResultCoverageOFoura}{\result{coverage}}
\newcommand{\ResultRateRatioHOne}{\result{ratio}}
\newcommand{\ResultRateRatioLOne}{\result{ratio}}
\newcommand{\ResultSlowShiftHOne}{\result{change}}
\newcommand{\ResultSlowShiftLOne}{\result{change}}
\newcommand{\ResultSlowFractionOThreeHOne}{\result{f}}
\newcommand{\ResultSlowFractionOThreeLOne}{\result{f}}
\newcommand{\ResultSlowFractionOFouraHOne}{\result{f}}
\newcommand{\ResultSlowFractionOFouraLOne}{\result{f}}

%  \generated{name}{fallback rows} inserts a pipeline-generated table body if
%  it exists and the placeholder rows otherwise, so the document compiles at
%  every stage of completeness.
%  \generated{name} pulls in the pipeline-written body rows of a table if
%  they exist, and the placeholder rows otherwise.  Note that the fallback is
%  a FILE, not an argument: table rows contain & and \\, and TeX cannot
%  absorb an alignment separator into a macro argument -- doing so silently
%  swallows every row into a single cell and the table then fails with a
%  misplaced \noalign at the \bottomrule.  Both branches here are bare
%  \input calls, so neither ever sees an & at all.
\newcommand{\generated}[1]{%
  \IfFileExists{generated/#1.tex}%
    {\input{generated/#1}}%
    {\input{placeholders/#1}}%
}

\IfFileExists{numbers.tex}{\input{numbers.tex}}{}

% ---- figures that may not exist yet ---------------------------------
%  \safefig[width]{name} includes results/figures/<name>.{pdf,png} if it is
%  there, and otherwise draws a labelled placeholder box.  This lets the
%  document compile before the production run has generated its figures,
%  without silently omitting them.
\makeatletter
\newcommand{\safefig}[2][\textwidth]{%
  \IfFileExists{../results/figures/#2.pdf}{\includegraphics[width=#1]{#2}}{%
  \IfFileExists{figures/#2.pdf}{\includegraphics[width=#1]{#2}}{%
  \IfFileExists{../results/figures/#2.png}{\includegraphics[width=#1]{#2}}{%
  \IfFileExists{figures/#2.png}{\includegraphics[width=#1]{#2}}{%
    \setlength{\fboxsep}{0pt}%
    \fcolorbox{black!25}{black!3}{%
      \parbox[c][0.26\textheight][c]{#1}{%
        \centering\small\ttfamily
        \textcolor{black!55}{figure not yet generated}\\[0.5em]
        \detokenize{#2}\\[0.5em]
        \textcolor{black!55}{\detokenize{python scripts/07_make_figures.py}}%
      }}%
  }}}}%
}
\makeatother

% ---- shorthands ------------------------------------------------------
\newcommand{\Hone}{\mbox{H1}}
\newcommand{\Lone}{\mbox{L1}}
\newcommand{\Otza}{\mbox{O3a}}
\newcommand{\Otzb}{\mbox{O3b}}
\newcommand{\Ofa}{\mbox{O4a}}
\newcommand{\Tarch}{T_{\mathrm{arch}}}
\newcommand{\fdrive}{f_{\mathrm{drive}}}
\newcommand{\code}[1]{\texttt{\small #1}}

% LIGO channel names are long, unbreakable, and appear inline throughout.
% Give TeX a little room rather than letting them push into the margin.
\setlength{\emergencystretch}{3.5em}
\hbadness=3000

% =====================================================================
\begin{document}

\begin{titlepage}
\centering
\vspace*{2.0cm}

{\Large\bfseries
Q-scan morphology of scattered-light transients\\[0.35em]
in LIGO O3 and O4a\par}

\vspace{0.9em}
{\large\itshape with an auxiliary-channel correlation analysis\par}

\vspace{2.4cm}
{\large Hardik Sher Singh\par}
\vspace{0.4em}
{\normalsize Department of Electronics and Communication Engineering\\
Thapar Institute of Engineering and Technology, Patiala, India\par}
\vspace{0.3em}
{\small\texttt{hsingh\_be23@thapar.edu}\par}

\vspace{2.8cm}
\begin{minipage}{0.85\textwidth}
\small
\noindent\textbf{Abstract.}
Scattered light is among the most persistent sources of transient noise in the
Advanced LIGO detectors. Stray light reflecting from a moving surface and
rejoining the main beam is phase modulated by $4\pi x/\lambda$, and when that
excursion exceeds a radian the noise appears in the time--frequency plane as a
train of harmonic \emph{arches} whose instantaneous frequency follows
$f_n = 2n|v|/\lambda$. Because the frequency depends on the magnitude of the
velocity, each cycle of the driving motion produces two arches, so the arch
repetition period is a direct measurement of the mechanical disturbance
responsible.

This work characterises the morphology of scattered-light transients in the
third and fourth LIGO observing runs using only public data. A labelled O3
sample is drawn from the Gravity Spy classification release; every trigger is
Q-transformed and reduced to a twenty-element feature vector describing the
arch train. The O3 sample is then tied to physical cause in two ways: by
rank-correlating the hourly trigger rate against seismometer band-limited RMS
in the earthquake, microseismic and anthropogenic bands, and by fitting the
$f_n = 2n|v|/\lambda$ predictor, trigger by trigger, to the Livingston OM3
suspension witness that the public O3 auxiliary release makes available at
\SI{256}{\hertz}.

Two properties of the public data determine the shape of the O4a half of the
analysis and are reported as findings in their own right: no Gravity Spy
classification catalogue has been released for O4a, and the O4a auxiliary
release contains no seismometer, accelerometer or suspension witness of any
kind. The O4a sample is therefore constructed from strain alone, using an
excess-power trigger finder and a classifier trained on the O3 features, and
the O3 comparison is made like-for-like by running the identical search on
both runs. The auxiliary analysis, being an O3 result, is what licenses
reading the O4a arch-period distribution as a statement about ground motion.

\vspace{0.6em}
\noindent\textbf{Keywords:} gravitational waves; LIGO; detector
characterisation; transient noise; scattered light; time--frequency analysis.
\end{minipage}

\vfill
{\small Analysis code and configuration:
\url{https://github.com/hardiksingh/LIGO_Glitch_Classification}\par}
\vspace{0.4em}
{\small \today\par}
\end{titlepage}


\tableofcontents
\clearpage

\section{Introduction}
\label{sec:introduction}

\subsection{Transient noise and why its shape matters}

The Advanced LIGO detectors \citep{AdvancedLIGO2015} reach strain sensitivities
of order $10^{-23}\,\mathrm{Hz}^{-1/2}$, and at that level the data are not
Gaussian. Short bursts of excess power --- \emph{glitches} --- occur at rates
of order one per minute, and they matter for two reasons: they reduce the
livetime available for astrophysical search, and they can mimic or corrupt real
signals \citep{Davis2021}. Detector characterisation exists to find them,
explain them, and where possible remove their cause from the instrument.

The productive question about a glitch is rarely ``how loud was it'' but
``what shape was it''. Glitch classes are defined morphologically, in the
time--frequency plane, and the morphology is what points back to a mechanism.
The Gravity Spy project \citep{Zevin2017,Glanzer2023} formalised this by
combining citizen-science labelling with a convolutional classifier over
$22$ morphological classes, and its classification release for the first three
observing runs is now the standard labelled dataset for this kind of work.

\subsection{Scattered light}

Among those classes, \emph{scattered light} is unusual in that its physical
mechanism is understood in closed form.

A small fraction of the circulating beam leaves the main path, strikes a
nearby surface --- a baffle, a suspended optic, a viewport --- and rejoins the
beam. Having traversed the extra path twice, the recombining field has
acquired a phase
%
\begin{equation}
  \phi(t) = \frac{4\pi\,x(t)}{\lambda},
  \label{eq:phase}
\end{equation}
%
where $x$ is the displacement of the scattering surface along the beam and
$\lambda = \SI{1064}{\nano\metre}$ is the Advanced LIGO laser wavelength. When
$\phi$ stays well below a radian the coupling is linear and the scattered light
simply adds a small amount of noise at the frequency of the motion. When $\phi$
exceeds a radian --- which requires only a fraction of a micron of
displacement --- the modulation is strongly non-linear, and the noise appears
at the instantaneous frequency
%
\begin{equation}
  f_1(t) = \frac{1}{2\pi}\left|\frac{\mathrm{d}\phi}{\mathrm{d}t}\right|
         = \frac{2}{\lambda}\,|v(t)|,
  \label{eq:f1}
\end{equation}
%
with $v = \mathrm{d}x/\mathrm{d}t$. Numerically,
$2/\lambda = \SI{1.88e6}{\hertz\per(\metre\per\second)}$: every micron per
second of surface velocity produces \SI{1.88}{\hertz} of noise frequency.

Light that makes $n$ round trips between the two surfaces before rejoining
accumulates $n$ times the phase, so the noise appears simultaneously at every
harmonic
%
\begin{equation}
  f_n(t) = n\,\frac{2}{\lambda}\,|v(t)|, \qquad n = 1, 2, 3, \ldots
  \label{eq:fn}
\end{equation}
%
with amplitudes falling off as the bounce order rises. In a Q-scan this
produces the characteristic picture: a comb of smooth, harmonically spaced
ridges, all peaking at the same instant.

Two consequences of \cref{eq:fn} carry the whole analysis.

\paragraph{The arch period measures the driving motion.} Because $f_n$ depends
on $|v|$ rather than on $v$, it reaches a maximum twice per cycle of the
motion. The observed arch repetition period is therefore \emph{half} the
period of the disturbance,
%
\begin{equation}
  \fdrive = \frac{1}{2\,\Tarch}.
  \label{eq:fdrive}
\end{equation}
%
A measured $\Tarch$ of \SI{3.3}{\second} implies a driver at
\SI{0.15}{\hertz}, squarely in the microseismic band; a measured
\SI{0.25}{\second} implies \SI{2}{\hertz}, in the anthropogenic band. The
morphology alone, with no auxiliary data at all, identifies the class of
disturbance.

\paragraph{The harmonic spacing measures the fundamental.} At the apex of an
arch train, all harmonics peak together, so a frequency slice through the apex
is a comb with spacing $f_1^{\max} = 2 v_{\mathrm{peak}}/\lambda$. The spacing
is thus a direct readout of the peak surface velocity, again from strain alone.

\subsection{What is already known}

\citet{Soni2021Scattering} characterised scattering throughout O3 and
identified two distinct populations at LIGO: a \emph{slow} component driven by
the \SIrange{0.1}{0.3}{\hertz} microseism, whose arches recur every few
seconds and reach a few tens of hertz, and a \emph{fast} component driven by
\SIrange{1}{6}{\hertz} motion, whose arches recur several times per second.
They traced the dominant O3 slow-scattering coupling to relative motion of a
few microns between the main suspension chain and its reaction chain, and
demonstrated mitigation by damping that relative motion. A witness-channel
algorithm predicting when scattering should be present is a central result of
that work, and \cref{eq:fn} is its basis.

Related work has established the surrounding picture: environmental coupling
paths in general \citep{Nguyen2021}, scattering from end benches at Virgo
\citep{Was2021,Accadia2010}, and the O3 data-quality products
\citep{Davis2021}. For the fourth run, \citet{Capote2025} document the
instrument performance and \citet{Soni2024O4} the detector characterisation of
its first half. Most directly relevant, \citet{Nandi2026} report on scattered
light at Livingston during O4 and again find two populations, one traced to
microseismic ground motion and mitigated by baffle installation, the other to
vertical ground motion coupling through a vacuum chamber.

That body of work is, however, built on internal LIGO data: full auxiliary
channel sets, engineering records, and access to the instrument. The present
study asks a narrower and quite different question.

\subsection{This work}

\emph{How much of the scattered-light story can be reconstructed from public
data alone, and does the answer change between O3 and O4a?}

The question is worth asking for two reasons. It is a genuine test of what the
open data releases support, which is useful to anyone outside the
collaboration. And it forces every inference to be made from quantities that
can be independently checked, which is a discipline that internal-data
analyses do not always impose.

The analysis proceeds in five steps.

\begin{enumerate}
  \item A labelled O3 sample of scattered-light triggers is drawn from the
        Gravity Spy release \citep{GravitySpyData2021}, together with a
        control sample of other glitch classes.
  \item Each trigger is Q-transformed \citep{Chatterji2004} and reduced to a
        twenty-element morphology feature vector, described in
        \cref{sec:methods}, whose components are the physical quantities of
        \crefrange{eq:fn}{eq:fdrive} rather than generic image statistics.
  \item The O3 sample is tied to physical cause using the public auxiliary
        channels: the hourly trigger rate is rank-correlated against
        seismometer band-limited RMS, and the predictor of \cref{eq:fn} is
        fitted trigger by trigger to a suspension witness.
  \item A classifier trained on the O3 features is applied to a
        livetime-limited blind search of O4a strain.
  \item The two populations are compared in rate, in slow/fast composition,
        and feature by feature.
\end{enumerate}

Steps 3 and 4 are shaped by two properties of the public releases that were
established before any code was written and that are reported here as results
in their own right, because they are not stated plainly anywhere else:

\begin{itemize}
  \item \textbf{No Gravity Spy classification catalogue exists for O4a.} The
        Zenodo release \citep{GravitySpyData2021} covers O1 through O3b. The
        O4a Gravity Spy model is described by \citet{Wu2025} but no catalogue
        accompanies it. Independent confirmation comes from
        \citet{Cirfeta2026}, who work on O4a without labels and use the O3b
        Gravity Spy set as their reference index --- exactly the workaround
        that O4a labels would render unnecessary.
  \item \textbf{The O4a auxiliary release contains no motion witness.} Its
        $25$ channels \citep{GWOSCO4Aux} are calibration lines, input
        mode-cleaner wavefront sensors, photon-calibrator readouts, mains
        monitors and parametric-instability monitors. There is no
        seismometer, accelerometer, microphone or suspension channel, so
        neither the rate correlation nor the per-trigger predictor can be
        computed for O4a from public data.
\end{itemize}

The O3 auxiliary release, by contrast, proved richer than expected:
$86$ full-rate channels and $402$ second-trend channels
\citep{GWOSCO3Aux}, including the complete ground-seismometer band-limited RMS
family across the earthquake, microseismic and anthropogenic bands at both
sites, and --- crucially --- \code{L1:SUS-OM3\_M1\_DAMP\_L\_IN1\_DQ} at
\SI{256}{\hertz}, a genuine optic-motion witness and the only full-rate one in
either release.

\Cref{sec:data} sets out the data and verifies its availability;
\cref{sec:methods} the analysis; \cref{sec:validation} the validation against
simulated injections with known parameters; \cref{sec:results} the results;
\cref{sec:discussion} their interpretation and limits.

\section{Data}
\label{sec:data}

Every dataset used here is public. Nothing in the analysis requires LIGO
credentials. All availability statements below were verified against the live
services on 2026-08-04 by the environment-check stage of the pipeline
(\cref{app:reproducibility}).

\subsection{Observing-run epochs}

\Cref{tab:epochs} lists the epochs analysed and their GPS boundaries, taken
from the GWOSC data-release pages.

\begin{table}[htbp]
\centering
\caption{Observing-run epochs. Strain from both detectors is public for all
three. O4b and O4c are excluded: their releases were partial at the time of
writing and they add nothing to the O3-versus-O4a comparison.}
\label{tab:epochs}
\small
\begin{tabular}{llrrl}
\toprule
Epoch & Run & GPS start & GPS end & UTC span \\
\midrule
\Otza & O3 & 1238166018 & 1253977218 & 2019-04-01 to 2019-10-01 \\
\Otzb & O3 & 1256655618 & 1269363618 & 2019-11-01 to 2020-03-27 \\
\Ofa  & O4 & 1368975618 & 1389456018 & 2023-05-24 to 2024-01-16 \\
\bottomrule
\end{tabular}
\end{table}

Strain is taken from the \SI{4}{\kilo\hertz} release, which is sufficient for
an analysis confined below \SI{256}{\hertz}. Analysis-ready segments come from
the GWOSC segment API using the \code{<IFO>\_DATA} flag; this is the correct
denominator for a glitch rate, since a rate computed against calendar time
would confound glitchiness with duty cycle.

\subsection{Gravity Spy classifications: O3 only}
\label{sec:data:gravityspy}

The labelled input is the Gravity Spy machine-learning classification release,
Zenodo record 5649212 \citep{GravitySpyData2021,Glanzer2023}. It contains every
Omicron \citep{Robinet2020} trigger with $\mathrm{SNR} > 7.5$ and peak
frequency in \SIrange{10}{2048}{\hertz}, together with the trigger metadata, a
$22$-element vector of class posteriors, and the arg-max label.

Selection cuts, applied identically to the target and control classes so that
both are drawn under the same conditions:

\begin{itemize}
  \item $\mathrm{SNR} \geq 7.5$ (the Gravity Spy trigger threshold);
  \item peak frequency in \SIrange{10}{2048}{\hertz};
  \item machine-learning confidence in the assigned class $\geq 0.90$;
  \item event time inside the epoch boundaries of \cref{tab:epochs}.
\end{itemize}

The confidence cut is the one real choice here. It trades sample size for
label purity, and at $0.90$ it removes the ambiguous tail without collapsing
the sample. \Cref{sec:results:systematics} reports how the results move when
it is varied.

The control sample consists of the classes \emph{Blip}, \emph{Koi\_Fish},
\emph{Low\_Frequency\_Burst}, \emph{Tomte}, \emph{Extremely\_Loud},
\emph{Repeating\_Blips}, \emph{Light\_Modulation}, \emph{Violin\_Mode},
\emph{Whistle}, \emph{Power\_Line} and \emph{No\_Glitch}. Its role is to define
what the classifier of \cref{sec:methods:classifier} must reject.

\paragraph{There is no O4a catalogue.} This is worth stating precisely because
the analysis design turns on it. The Zenodo ML-classification release covers
O1, O2, O3a and O3b. The volunteer-classification release extending to
2024-07-20 is keyed by Zooniverse subject identifier and carries no GPS time,
detector or trigger metadata, so it cannot serve as a trigger list. The O4a
Gravity Spy model of \citet{Wu2025} is published without an accompanying
catalogue. Consequently the O4a sample must be built from strain
(\cref{sec:methods:search}).

\subsection{Auxiliary channels}
\label{sec:data:aux}

Both GWOSC auxiliary releases are served over NDS2 at \code{nds.gwosc.org}.
Only short windows around trigger times, and hour-scale trends, are needed
here, so the several-terabyte bulk downloads are unnecessary.

\subsubsection{O3}

The O3 release \citep{GWOSCO3Aux} comprises $86$ full-rate channels
($40$ \Hone, $46$ \Lone) and $402$ second-trend channels at \SI{1}{\hertz}
($190$ \Hone, $212$ \Lone). \Cref{tab:auxchannels} lists those the analysis
uses.

\begin{table}[htbp]
\centering
\caption{Auxiliary channels used in the O3 analysis. Names ending
\code{.rms} come from the second-trend frames. The suffix
\code{BLRMS\_100M\_300M} denotes a band-limited RMS over
\SIrange{0.1}{0.3}{\hertz}; ``M'' is millihertz.}
\label{tab:auxchannels}
\footnotesize
\begin{tabular}{llp{6.4cm}}
\toprule
Channel & Rate & Role \\
\midrule
\code{L1:SUS-OM3\_M1\_DAMP\_L\_IN1\_DQ}
  & \SI{256}{\hertz}
  & OM3 top-stage \emph{length} damping error signal; the motion witness for
    the per-trigger predictor \\[0.3em]
\code{[HL]1:ISI-GND\_STS\_*\_BLRMS\_100M\_300M.rms}
  & \SI{1}{\hertz}
  & microseismic band at ETMX, ETMY, ITMX, ITMY, HAM5 in $x$, $y$, $z$;
    driver of slow scattering \\[0.3em]
\code{[HL]1:ISI-GND\_STS\_*\_BLRMS\_\{1\_3,3\_10\}.rms}
  & \SI{1}{\hertz}
  & \SIrange{1}{3}{\hertz} and \SIrange{3}{10}{\hertz} anthropogenic bands;
    driver of fast scattering \\[0.3em]
\code{[HL]1:ISI-GND\_STS\_ITMY\_Z\_BLRMS\_30M\_100M}
  & \SI{16}{\hertz}
  & \SIrange{0.03}{0.10}{\hertz} earthquake band \\
\bottomrule
\end{tabular}
\end{table}

Two asymmetries are worth recording. First, \Lone{} has a full-rate suspension
witness in the public release and \Hone{} does not; the per-trigger predictor
of \cref{sec:methods:predictor} is therefore an \Lone-only result, and the
\Hone{} auxiliary analysis is confined to the rate level. Second, the
auxiliary channels are published in raw digital counts with no calibration to
metres, which is dealt with in \cref{sec:methods:predictor}.

\subsubsection{O4a}
\label{sec:data:aux:o4a}

The O4a release \citep{GWOSCO4Aux} is $25$ channels: $14$ from \Hone{} and
$11$ from \Lone. They are reproduced in full below, because the point being
made depends on the list being complete rather than summarised.

\begin{lstlisting}[language=,basicstyle=\ttfamily\scriptsize]
H1: CAL-PCALX_RX_PD_OUT_DQ            L1: ASC-X_TR_A_NSUM_OUT_DQ
    CAL-PCALY_RX_PD_OUT_DQ                CAL-PCALX_RX_PD_OUT_DQ
    IMC-WFS_A_DC_PIT_OUT_DQ               CAL-PCALY_RX_PD_OUT_DQ
    IMC-WFS_A_DC_YAW_OUT_DQ               IMC-WFS_B_I_PIT_OUT_DQ
    LSC-POP_A_LF_OUT_DQ                   OAF-IMC_WFS_B_I_PIT_PREFILT_OUT_DQ
    LSC-REFL_A_RIN_OUT_DQ                 PEM-EX_VMON_ETMX_ESDPOWER24_DQ
    OAF-IMC_WFS_A_DC_PIT_PREFILT_OUT_DQ   PEM-EY_MAINSMON_EBAY_1_DQ
    OAF-IMC_WFS_A_DC_YAW_PREFILT_OUT_DQ   SUS-ETMX_L1_CAL_LINE_OUT_DQ
    OAF-REFL_A_RIN_PREFILT_OUT_DQ         SUS-ETMX_L2_CAL_LINE_OUT_DQ
    PEM-EY_MAINSMON_EBAY_1_DQ             SUS-ETMX_L3_CAL_LINE_OUT_DQ
    SUS-ETMX_L1_CAL_LINE_OUT_DQ           SUS-PI_PROC_COMPUTE_MODE5_RMSMON
    SUS-ETMX_L2_CAL_LINE_OUT_DQ
    SUS-ETMX_L3_CAL_LINE_OUT_DQ
    SUS-PI_PROC_COMPUTE_MODE29_RMSMON
\end{lstlisting}

Every one is a photon-calibrator readout, an input mode-cleaner wavefront
sensor, a length- or angular-sensing readout, a mains voltage monitor, an
electrostatic-drive voltage monitor, a calibration-line drive, or a
parametric-instability RMS monitor.

There is no seismometer, no accelerometer, no microphone and no suspension
witness. Neither the rate-level correlation against ground motion nor the
per-trigger scattering predictor is computable for O4a from public data. This
is a limitation of the release, not of the method, and it is the reason the
O3-versus-O4a comparison in \cref{sec:results:comparison} is a morphology
comparison rather than a mechanism comparison.

The analysis pins this fact in code: a unit test asserts that the O4a channel
list contains none of the tokens \code{ISI-GND}, \code{\_SEIS\_},
\code{\_ACC\_}, \code{SUS-OM} or \code{WIT\_L}. Should GWOSC extend the
release, that test fails, which is the intended signal to revisit the design.

\subsection{Summary of what is and is not possible}

\begin{table}[htbp]
\centering
\caption{Analysis coverage by run, given the public data.}
\label{tab:coverage}
\small
\begin{tabular}{lll}
\toprule
Analysis & O3 & \Ofa \\
\midrule
Q-scan morphology            & yes, on Gravity Spy labels & yes, on classifier output \\
Rate versus ground motion    & yes                        & \textbf{no} --- no seismic channels \\
Per-trigger predictor        & \Lone{} only               & \textbf{no} --- no motion witness \\
Slow/fast decomposition      & yes                        & yes, from arch period alone \\
\bottomrule
\end{tabular}
\end{table}

The slow/fast decomposition survives into O4a precisely because
\cref{eq:fdrive} makes it a strain-only measurement. That is what makes the
comparison possible at all, and it is the reason the auxiliary work in
\cref{sec:results:aux} matters beyond O3: it is the calibration that justifies
reading an O4a arch period as a statement about ground motion.

\section{Method}
\label{sec:methods}

\subsection{Q-transform maps}
\label{sec:methods:qscan}

Each trigger is analysed on a normalised-energy Q-transform map
\citep{Chatterji2004}. The transform tiles the time--frequency plane with
constant-$Q$ bisquare windows and reports, for each tile,
%
\begin{equation}
  Z = \frac{|X|^2}{\langle |X|^2 \rangle},
  \label{eq:normenergy}
\end{equation}
%
the ratio of tile energy to the mean over tiles at the same frequency. For
stationary Gaussian noise $Z$ is exponentially distributed with unit mean, so
$\sqrt{2Z}$ is the single-tile signal-to-noise ratio. The implementation is
\code{gwpy}'s \code{q\_transform} \citep{Macleod2021GWpy}, which is the same
algorithm used by the LIGO detector-characterisation tools; the standard
\code{gwdetchar} wrapper is not used because its dependency chain has no
Windows build, and nothing beyond the transform itself is needed.

\SI{64}{\second} of strain is fetched around each trigger and only the central
\SI{16}{\second} is analysed. The padding is not decorative: the whitening
filter and the constant-$Q$ filters both ring, and the resulting transients
are exactly the kind of structure a ridge tracker would mistake for an arch.
The configuration enforces that the padding exceeds the analysed window.

Two parameters were set by validation rather than by convention, and both
mattered.

\paragraph{The $Q$ range is capped at 16.} \code{gwpy} selects the $Q$ plane of
maximum energy. Left unconstrained it selects $Q \approx 45$ for a scattering
train, because a high-$Q$ plane integrates the quasi-periodic structure
coherently and reports enormous energy. But a tile of quality factor $Q$ at
frequency $f$ lasts $Q/2\pi f$ seconds --- \SI{0.16}{\second} at $Q = 45$ and
\SI{45}{\hertz} --- which is comparable to the repetition period of fast
scattering. The arches are then unresolved. On a simulated fast event with an
injected arch period of \SI{0.250}{\second}, the uncapped transform returned
\SI{4.99}{\second}: wrong by a factor of twenty. With the cap it returns
\SI{0.250}{\second}.

\paragraph{The ridge threshold is $Z = 20$, not $Z = 5.5$.} A threshold has to
sit above the largest energy a \emph{noise-only} map produces, not merely
above its mean. A map of this size holds of order $10^4$ independent tiles,
whose expected maximum is $\ln(10^4) \approx 9$ before the multi-$Q$ search is
taken into account and $\approx 16$ after. At $Z = 5.5$ --- a natural-looking
choice, corresponding to single-tile SNR $3.3$ --- a pure-noise map produced
hundreds of supra-threshold bins and the measured maximum arch frequency
saturated at the top of the analysis band. $Z = 20$ (single-tile SNR $6.3$)
clears the noise; a connected-component filter removes the residue.

\subsection{Ridge extraction}
\label{sec:methods:ridge}

From the thresholded map, three tracks are formed per time bin: the frequency
of the loudest tile, the highest supra-threshold frequency, and the lowest.
The second of these --- the \emph{arch envelope} --- is the quantity the
auxiliary predictor predicts, and is what the morphology features describe.

Connected supra-threshold regions smaller than $40$ output pixels are
discarded. Because the output grid is interpolated from far coarser $Q$ tiles,
a single hot tile already spans tens of pixels, so this is a speck filter
rather than a significance cut; the significance is carried by the threshold.

\subsection{Morphology features}
\label{sec:methods:features}

Each map is reduced to the twenty features listed in \cref{app:features}.
They fall into five groups: extent (peak, maximum and minimum frequency,
bandwidth, duration, SNR); arch structure (count, repetition period, period
scatter); harmonic structure (count and spacing at the apex); shape (RMS slope
and curvature of the envelope, ratio of rising to falling time); and energy
distribution (fraction below \SI{40}{\hertz}, fraction above
\SI{100}{\hertz}, spectral centroid, band concentration, time symmetry).

They are computed on the normalised-energy map, so they are insensitive to the
absolute calibration and directly comparable between detectors and between
runs. Where a feature is undefined --- \code{arch\_period} when fewer than two
arches are found --- it is returned as \code{NaN} rather than as a fabricated
value, and the classifier treats the missingness as informative.

Two of them deserve comment.

\paragraph{Arch period, from the autocorrelation.} The obvious estimator is
the median spacing between detected envelope maxima. It fails in exactly the
case that matters: when the topmost harmonic dips below threshold the envelope
drops to the harmonic below, creating a spurious local maximum that splits one
arch into two. On a simulated slow event with an injected period of
\SI{3.33}{\second}, peak spacing returned \SI{2.80}{\second} --- \SI{16}{\percent}
low. The autocorrelation of the envelope is insensitive to this, because the
periodicity is present in every harmonic simultaneously; it returns
\SI{3.36}{\second}. The peak count is still reported separately, as
\code{n\_arches}.

\paragraph{Maximum arch frequency is biased high, and stays that way.} A $Q$
tile has bandwidth $f/Q$, so the supra-threshold region extends above the true
apex of the arch by roughly half of that. Measured against injections, the
bias is about \SI{15}{\percent}. It is not tuned away: it is measured,
asserted in the test suite, and quoted with the result. Since it applies
identically to both runs it cancels in the O3-versus-O4a comparison, which is
where the feature is actually used.

\subsection{From triggers to events, and from events to blocks}
\label{sec:methods:sampling}

Two properties of the real data, neither visible until the catalogue and the
strain server were in hand, determine how the sample is built. Both are
recorded here because each one changes a number that would otherwise be
reported.

\paragraph{The catalogue lists triggers, not events.} Gravity Spy publishes
Omicron triggers. After the cuts of \cref{sec:data:gravityspy},
\SI{12.3}{\percent} of rows share an exact $(\mathrm{event\ time},
\mathrm{detector})$ pair -- the same trigger under several Gravity Spy subject
identifiers, one per Q-scan duration view -- and about half of all triggers
fall within \SI{8}{\second} of another one, which is to say inside each other's
analysed Q-scan window. Those triggers therefore produce \emph{identical}
feature vectors.

Left alone, every morphology distribution would be weighted by how many
Omicron triggers a given arch train happened to generate: loud, long trains
counted dozens of times, quiet ones once. That is a selection effect acting
directly on the axis this study measures. Worse, it is class-dependent.
Clustering at \SI{16}{\second} -- the full width of the analysed window, so
that no two retained events can share a Q-scan -- yields \num{2.44} triggers
per event for scattered light but only \num{1.35} for the control classes.
Uncorrected, scattering would have been over-represented by some
\SI{80}{\percent} relative to the controls, biasing both the classifier's
training set and every rate comparison. The same window is applied to the O4a
search output, so the two runs count the same kind of thing.

The unclustered trigger list is retained as well: the rate correlation of
\cref{sec:methods:ratecorr} legitimately uses \emph{trigger} rate, and each
result states which of the two it quotes.

\paragraph{The unit of cost is the file, not the event.} GWOSC serves open
strain in HDF5 files of \SI{4096}{\second}, and a request for any
sub-interval downloads the whole containing file: two to four minutes on a
domestic connection, essentially independent of how much was asked for. A
\SI{64}{\second} Q-scan window therefore costs one full download, and a
time-stratified sample of \num{4800} individual events touched \num{3484}
distinct files -- about \num{1.4} events per download, or three days of wall
clock.

The analysis is restructured around that. Windows are cut from whole blocks
held in a small in-memory cache, and the sample is drawn as \emph{blocks
first, events within blocks second}: \num{100} blocks per detector per epoch,
stratified across each epoch's span, with up to \num{12} events taken at
random from each. That gives \num{2256} scattered-light events from \num{400}
blocks, and the measured yield in a production run is about ten events per
download. The control sample is drawn from the \emph{same} blocks, which costs
no extra downloads and removes a confound at the same time: signal and
background then come from identical stretches of data, so the classifier
cannot separate the two classes by learning what a noisy night looks like.

The price is real and is carried into the interpretation. Events drawn from
one block share ground-motion conditions and are therefore correlated, so the
effective sample size is smaller than the row count: it is the \emph{block}
count, not the event count, that carries the statistical weight. Events are
taken at random within a block rather than by loudness, which would bias the
sample towards the tail.

\subsection{Rate-level auxiliary correlation}
\label{sec:methods:ratecorr}

Triggers are binned into GPS hours and divided by the analysis-ready livetime
in each hour. Hours with zero livetime and hours with zero triggers are
different things: the former must be dropped, the latter kept as genuine
zeros. Conflating them is a reliable way to manufacture a spurious correlation
with anything that tracks duty cycle, ground motion included.

Each hourly rate series is rank-correlated (Spearman $\rho$) against the
hourly median band-limited RMS of every configured seismometer channel. The
median within an hour is used rather than the mean because the trend channels
contain occasional large excursions --- earthquakes, calibration activity ---
whose effect on an hourly mean is out of proportion to how much of the hour
they occupy.

Fetching those trends has its own constraints, measured against the live
server rather than assumed. The GWOSC NDS host serves a four-hour request but
refuses eight; nine channels requested together cost \SI{11.3}{\second}
against \SI{40.8}{\second} one at a time, so requests are batched; and
requests are issued only for chunks that overlap analysis-ready segments,
since a chunk with no livetime contributes nothing to a rate correlation.
Availability looked poor when probed at fixed calendar offsets and turned out
to be complete -- 18 of 18 -- when probed uniformly in livetime, the earlier
failures having been non-observing time.

Confidence intervals come from a moving-block bootstrap
\citep{Kunsch1989} with a six-hour block. This matters: the microseism varies
on a timescale of days, so hourly rates are strongly autocorrelated and the
naive standard error on $\rho$ is far too small. Three further checks guard
against the standard failure modes:

\begin{itemize}
  \item a \textbf{quantile table}, giving the mean rate in each quartile of
        the driving channel and the top-to-bottom ratio, so that the result is
        stated in units a commissioner can act on rather than as a correlation
        coefficient;
  \item a \textbf{lag scan} over $\pm\SI{12}{\hour}$. Scattering responds to
        ground motion essentially instantaneously, so the correlation must
        peak at zero lag. A peak at large lag would indicate that both series
        are being driven by a slow common trend --- diurnal or weather ---
        rather than by the mechanism under test;
  \item a \textbf{partial correlation}, testing whether an apparent
        correlation with the \SIrange{1}{6}{\hertz} anthropogenic band
        survives once the microseism is controlled for.
\end{itemize}

\subsection{Trigger-level scattering predictor}
\label{sec:methods:predictor}

For \Lone, where a full-rate motion witness is public, the predictor of
\cref{eq:fn} is tested trigger by trigger.

The witness is band-passed over \SIrange{0.03}{10}{\hertz} --- the lower
corner removes the DC drift that differentiation would otherwise amplify into
a ramp --- and differentiated to give velocity. Harmonics $n = 1 \ldots 8$ of
\cref{eq:fn} are then compared with the measured arch envelope. A ridge point
at frequency $f$ counts as \emph{explained} if some harmonic lies within
$\max(\SI{1.5}{\hertz},\, 0.15 f)$ of it, and the headline statistic is the
fraction of supra-threshold time bins that are explained.

\paragraph{The one free parameter.} GWOSC publishes auxiliary channels in raw
counts with no calibration to metres. Rather than invent one, a single scale
factor $\alpha$ (metres per count) is fitted per detector and epoch by
maximising the mean explained fraction over a training set of triggers, using
a logarithmic grid over eight decades followed by a golden-section refinement.
It is then \emph{frozen} and applied to a disjoint validation set, on which the
quoted agreement is measured. The width of the range over which the objective
stays within \SI{1}{\percent} of its maximum is reported alongside: a narrow
range means the fit is well constrained, a wide one is a warning that the
witness does not in fact track the scattering surface.

\Cref{sec:validation} shows this procedure recovering a known calibration from
a synthetic witness to \SI{7}{\percent}.

\subsection{The strain-only identifier}
\label{sec:methods:classifier}

To carry the O3 definition of the class into O4a, a gradient-boosted tree
classifier \citep{Friedman2001,Pedregosa2011} is trained on the O3 features to
separate \emph{Scattered\_Light} from the control classes. Missing features are
median-imputed with a missingness indicator, so that ``no arch period could be
measured'' is available to the model as information rather than being silently
filled in.

Performance is measured by stratified five-fold cross-validation, and the
operating point is fixed at a false-positive rate of $0.01$ against the
control sample, then frozen.

A third constraint is arithmetic and binds harder than either of the caveats
below. An operating point can only be placed as finely as the background
sample resolves: at a false-positive rate $p$ with $n$ control events per fold,
the threshold is fixed by $p\,n$ expected false positives. Trained on a trial
sample of 66 controls, the cross-validated AUC was a stable
$0.919 \pm 0.033$ while the efficiency at $p = 0.01$ came out
$0.368 \pm 0.334$, with individual folds returning $0.000$ (at an infinite
threshold), $0.160$, $0.680$, $0.760$ and $0.240$ --- because $p\,n = 0.13$,
so the threshold was decided by whether zero or one event happened to land
above it. The code now computes $p\,n$, warns below five, and records the fact
in the saved model; the production control sample gives $p\,n \approx 7$ and
clears it. The AUC is threshold-free and unaffected throughout.

Two further caveats travel with those numbers and are repeated wherever they
are quoted. The efficiency and false-positive rate are measured against the
\emph{Gravity Spy control classes}, so they say how well the model separates
scattered light from other known O3 glitch classes; they do not and cannot
measure performance against whatever unlabelled O4a contains. And the folds
are not grouped in time, so triggers minutes apart can land in different
folds, making the quoted AUC mildly optimistic. The sample is drawn stratified
across the whole run, which makes such neighbours rare but does not eliminate
them.

\subsection{The O4a search}
\label{sec:methods:search}

With no O4a trigger catalogue, triggers are generated from strain. The
statistic is the \SI{0.25}{\second} RMS of whitened, band-passed
(\SIrange{10}{120}{\hertz}) strain, converted to a robust $z$-score using the
median and median-absolute-deviation of the surrounding stretch. Robust
statistics are essential here: during a scattering-heavy hour a large fraction
of the data is glitching, and a mean/standard-deviation normalisation would be
dragged upward by the very signal it is meant to detect. Candidates are
clustered with a \SI{4}{\second} window, since an arch train is one physical
event lasting seconds, and candidates within the whitening-filter transient at
either end are discarded.

This is a deliberately loose stand-in for Omicron: its job is to lose no
scattering, not to be pure. Purity is the classifier's job.

\paragraph{Livetime, not calendar time.} A full O4a search is out of reach on
the hardware available. The search draws $1200$ windows of \SI{512}{\second}
per detector, sampled uniformly in livetime --- about \SI{7.1}{\day} each,
roughly \SI{4}{\percent} of the run. All rates are therefore quoted per hour of
\emph{analysed} livetime, and the coverage fraction is recorded with every
result.

\paragraph{Like-for-like comparison.} The identical search is also run on O3b.
This is what makes the comparison interpretable: with the same trigger
generator and the same classifier on both sides, a difference in rate or
morphology cannot be attributed to the change from the Gravity Spy pipeline to
this one. The comparison against the Gravity Spy O3 sample is also reported,
since it has better labels; where the two agree the result is robust, and
where they disagree the difference is a selection effect and is reported as
such.

\subsection{Population statistics}
\label{sec:methods:stats}

Rates are Poisson counts over a known livetime, so intervals are exact
(Garwood) Poisson intervals rather than $\sqrt{N}$.

Feature distributions are compared with the two-sample Kolmogorov--Smirnov
test for the $p$-value and with Cliff's $\delta$ \citep{Cliff1993} for the
effect size. With tens of thousands of triggers per sample essentially every
KS test rejects, so the $p$-value carries little information; the effect size
is what carries the physics and is quoted first. Because many features are
tested at once, $p$-values are corrected by the Benjamini--Hochberg procedure
\citep{BenjaminiHochberg1995} and the corrected values are the ones reported.

The headline slow-fraction change is a difference of two binomial proportions,
reported with a Newcombe score-based interval \citep{Newcombe1998}, which
behaves correctly as a proportion approaches $0$ or $1$ where the Wald
interval does not.

\section{Validation against known signals}
\label{sec:validation}

Every number in \cref{sec:results} is produced by an algorithm that could, in
principle, be measuring an artefact of its own construction. The only way to
tell is to give it a signal whose parameters are known exactly and check what
comes back. This section does that, and it is the reason two of the analysis
choices in \cref{sec:methods} are what they are.

\subsection{The injection model}

The simulator generates
%
\begin{align}
  x(t) &= X_0 \sin(2\pi f_m t + \varphi), \label{eq:injx}\\
  h(t) &= \sum_{n=1}^{N} K\,r^{\,n-1}
          \sin\!\left(n\,\frac{4\pi\,x(t)}{\lambda}\right), \label{eq:injh}
\end{align}
%
where the sum runs over bounce orders with amplitude ratio $r$ between
successive bounces. The arches are therefore not drawn on: they emerge from
the phase modulation, as they do in the instrument. From
\crefrange{eq:injx}{eq:injh} the expected measurements follow:
%
\begin{equation}
  f_1^{\max} = \frac{4\pi f_m X_0}{\lambda}, \qquad
  f^{\max} = N f_1^{\max}, \qquad
  \Tarch = \frac{1}{2 f_m},
  \label{eq:injtruth}
\end{equation}
%
and the harmonic comb at the apex has spacing $f_1^{\max}$.

\paragraph{A wrong model, and what it cost.} The first version of the
simulator used a single term, $h = K \sin(4\pi x/\lambda)$, on the reasoning
that the Jacobi--Anger expansion of a sine of a large argument produces
harmonics automatically. It does --- but they are harmonics of the
\emph{motion} frequency $f_m$, and they terminate at $k \approx \beta$ where
$\beta = 4\pi X_0/\lambda$ is the modulation index. Since $\beta f_m$ is
exactly $f_1^{\max}$, a one-bounce model produces \emph{one} arch however
large $\beta$ is. Raising the displacement to ``get more harmonics'' only
makes that single arch taller. The harmonic arches of a real Q-scan come from
multiple round trips, which is what \cref{eq:injh} encodes.

A second false start is worth recording because it recurs in real analyses. A
scattering train that fills the entire analysed window is whitened away
completely: the transform estimates its amplitude spectral density from the
data it is given, and a signal occupying the whole stretch \emph{is} the noise
floor as far as that estimate is concerned. The injection is therefore
confined to a Tukey-windowed stretch occupying a modest fraction of the
window, as a real transient is.

\subsection{Recovery}

Two canonical injections are used throughout: a \emph{slow} event
($X_0 = \SI{8}{\micro\metre}$, $f_m = \SI{0.15}{\hertz}$, four bounces) and a
\emph{fast} event ($X_0 = \SI{0.64}{\micro\metre}$, $f_m = \SI{2}{\hertz}$,
three bounces), both at unit-variance white noise. \Cref{tab:validation} gives
the recovery.

\begin{table}[htbp]
\centering
\caption{Recovery of injected parameters. The witness calibration is fitted by
the procedure of \cref{sec:methods:predictor} without being told the answer;
the synthetic witness was generated at $10^{9}$ counts per metre, so the value
to be recovered is $1.00\times10^{-9}\,\mathrm{m}\,\mathrm{count}^{-1}$.}
\label{tab:validation}
\small
\begin{tabular}{lcccc}
\toprule
& \multicolumn{2}{c}{Slow injection} & \multicolumn{2}{c}{Fast injection} \\
\cmidrule(lr){2-3}\cmidrule(lr){4-5}
Quantity & injected & recovered & injected & recovered \\
\midrule
Arch period $\Tarch$ (\si{\second})   & 3.33 & 3.36 & 0.250 & 0.250 \\
Harmonic spacing (\si{\hertz})        & 14.2 & 13.9 & 15.1  & 14.0 \\
Arch count                            & $\sim5$ & 6 & $\sim32$ & 31 \\
Maximum arch frequency (\si{\hertz})  & 56.7 & 65.8 & 45.4 & 62.8 \\
Witness calibration ($10^{-9}$ m/count) & 1.00 & 1.07 & 1.00 & 1.09 \\
Envelope explained by predictor       & --- & \SI{95}{\percent} & --- & \SI{84}{\percent} \\
Predictor--envelope correlation $r$   & --- & 0.92 & --- & 0.83 \\
\bottomrule
\end{tabular}
\end{table}

The arch period, which carries the physical interpretation through
\cref{eq:fdrive}, is recovered to \SI{0.8}{\percent} for the slow injection and
exactly for the fast one. The harmonic spacing, which measures the peak
surface velocity, is recovered to within \SI{7}{\percent}. The witness
calibration --- fitted blind, over an eight-decade search range --- comes back
within \SI{9}{\percent}.

The maximum arch frequency reads systematically high, by \SI{16}{\percent} for
the slow injection and \SI{38}{\percent} for the fast one, for the reason given
in \cref{sec:methods:features}: a $Q$ tile of bandwidth $f/Q$ extends above
the true apex. The bias is larger for the fast injection because its arches
are shorter-lived relative to the tile duration. This feature is therefore not
used for any absolute statement; it is used only in the O3-versus-O4a
comparison, where the bias applies identically on both sides.

\begin{figure}[htbp]
\centering
\safefig[\textwidth]{validation_injections}
\caption{Q-scans of the two canonical injections with the extracted arch
envelope overlaid (white). Upper: slow scattering, \SI{0.15}{\hertz} drive,
arches recurring every \SI{3.3}{\second}. Lower: fast scattering,
\SI{2}{\hertz} drive, arches recurring every \SI{0.25}{\second}. Colour is
normalised energy $Z$; the ridge threshold is $Z = 20$. Produced by
\code{scripts/07\_make\_figures.py --only validation}.}
\label{fig:validation}
\end{figure}

\subsection{The null test}

A threshold that lets noise through invalidates every extent feature at once,
so the complementary test is as important as the recovery test: on a
pure-noise map at the configured threshold, the extractor must find
\emph{nothing}. It does --- zero supra-threshold time bins, zero duration, and
\code{NaN} for every frequency extent --- with the peak noise energy at
$Z = 14.6$, comfortably below the threshold of $20$. At the initially chosen
$Z = 5.5$ the same map produced hundreds of supra-threshold bins.

\subsection{A closure test on real triggers}
\label{sec:validation:closure}

Injections check the extractor against a signal built to the same model it
assumes, which is necessary but circular. A weaker but non-circular check is
available on real data: the features are over-determined, so they can be asked
whether they agree with each other.

On the first \num{126} scattered-light events extracted from \Hone{} \Otza,
the median measurements are an arch repetition period of \SI{0.21}{\second}, a
harmonic comb spacing of \SI{15.0}{\hertz}, and a maximum arch frequency of
\SI{41.5}{\hertz}. Those three are measured independently -- from the
autocorrelation of the envelope, from a frequency slice through the apex, and
from the extent of the supra-threshold region. Taken together they imply

\begin{itemize}
  \item a driving frequency $\fdrive = 1/2\Tarch = \SI{2.4}{\hertz}$, which
        falls in the anthropogenic band;
  \item a peak surface velocity $v = \lambda f_1^{\max} / 2 =
        \SI{8.0}{\micro\metre\per\second}$;
  \item a displacement amplitude $X_0 = \lambda f_1^{\max} / 4\pi \fdrive =
        \SI{530}{\nano\metre}$;
  \item $f^{\max}/f_1^{\max} = 2.8$ resolved harmonics.
\end{itemize}

Half a micron of optic motion at a few hertz, coupling through two or three
bounce orders, is an unremarkable description of fast scattering --- which is
the point. Had the extractor been measuring an artefact of its own thresholds,
there would be no reason for these three independent numbers to combine into a
physically sensible displacement.

The arch period is measurable for \SI{56}{\percent} of events and the harmonic
spacing for \SI{29}{\percent}; the remainder are too short or too faint for the
autocorrelation and the comb finder respectively. That completeness is a
property of the sample, not a failure, but it is why the classifier is built to
treat a missing feature as informative rather than imputing it silently.

\subsection{What only the real data revealed}
\label{sec:validation:realdata}

Injections test the physics; they cannot test the archive. Three defects
survived the entire simulated-signal suite and appeared within minutes of the
first contact with live data. They are recorded because each one is a class of
mistake that a validation-against-injections strategy is structurally blind
to.

\paragraph{The cache never wrote anything.} Both cache writers built a
temporary file, wrote it, and renamed it into place. But
\code{numpy.savez\_compressed} appends \code{.npz} to any filename that does
not already end in it, so the data landed in \code{<hash>.npz.tmp.npz} while
the rename looked for \code{<hash>.npz.tmp}. The very first real fetch died on
it. No injection test touched the cache, so nothing caught it; the suite now
round-trips both writers and checks that a failed write leaves nothing behind.

\paragraph{The catalogue counts triggers, and unevenly.} Described in
\cref{sec:methods:sampling}: the class-dependence is the part that matters,
since \num{2.44} triggers per scattered-light event against \num{1.35} per
control event is an \SI{80}{\percent} relative over-representation of the very
class under study. Nothing about a simulated arch train would have exposed
this, because the simulator produces one event per event.

\paragraph{A fifth of the trial sample was in a hole.} Twelve of the first
forty events failed, all inside one \SI{4096}{\second} block that GWOSC does
not serve --- neither as a whole file nor as a narrow window within it. The
error handling was already correct, in that the stage logged the failures and
carried on; what was wrong was the cost, since every event in the dead block
spent the full retry budget twice over. Unavailable blocks are now remembered,
and the count of events lost to data gaps is reported with each stage rather
than being silently absent from the sample.

Taken together these justify the shape of the suite: injections fix the
physics and the estimator, and the first real run is treated as a test in its
own right rather than as the start of production.

\subsection{Test suite}

The results above are encoded as assertions rather than left as a one-off
check: $114$ tests cover the feature contract, the recovery tolerances of
\cref{tab:validation}, the null test, the calibration round trip, a
negative control in which a predictor wrong by three decades must not still
appear to agree, the GPS-time conversions against the published epoch
boundaries, the segment arithmetic, the statistical estimators, and the two
data-availability facts of \cref{sec:data}. They require no network access and
run in about three minutes.

\section{Results}
\label{sec:results}

\IfFileExists{numbers.tex}{}{%
\begin{center}
\fcolorbox{red!60!black}{red!5}{%
\begin{minipage}{0.94\textwidth}
\small\textbf{Status.} The pipeline has been validated end to end against
simulated injections (\cref{sec:validation}) and every data path has been
verified against the live services, but the production run over the full O3
and O4a samples has not yet been executed. Quantities awaiting that run are
shown as \result{like this}. They are placeholders, not measurements, and no
value in this section should be quoted until this box disappears --- it does
so automatically once \code{paper/numbers.tex} is written by stage~06 of the
pipeline. The structure below, the statistical treatment, and the caveats are
final.
\end{minipage}}
\end{center}
\vspace{0.6em}}

\subsection{The O3 sample}
\label{sec:results:sample}

\Cref{tab:sample} gives the size of the labelled sample after the cuts of
\cref{sec:data:gravityspy}.

\begin{table}[htbp]
\centering
\caption{Gravity Spy triggers surviving the selection cuts, by detector and
epoch, and the scattered-light fraction. Produced by
\code{results/tables/catalog\_summary.csv}.}
\label{tab:sample}
\small
\generated{tab_sample}
\end{table}

\begin{figure}[htbp]
\centering
\safefig[0.9\textwidth]{example_slow_1_L1}
\caption{A representative slow-scattering trigger with the extracted arch
envelope overlaid. The harmonic comb, the smooth rise and fall, and the
regular repetition are the features the analysis quantifies.}
\label{fig:exampleslow}
\end{figure}

\subsection{Morphology of the O3 population}
\label{sec:results:morphology}

The distribution of the arch repetition period is the central descriptive
result, because \cref{eq:fdrive} converts it directly into a statement about
the mechanical disturbance. \Cref{fig:archperiod} shows it with the implied
driving frequency on the upper axis.

\begin{figure}[htbp]
\centering
\safefig[0.9\textwidth]{arch_period_o3_vs_o4a}
\caption{Arch repetition period. The upper axis gives the implied driving
frequency $\fdrive = 1/2\Tarch$; the microseismic band
(\SIrange{0.1}{0.3}{\hertz}) and the anthropogenic band
(\SIrange{1}{6}{\hertz}) can be read off directly. The dotted line at
\SI{1.5}{\second} is the slow/fast split.}
\label{fig:archperiod}
\end{figure}

\Cref{tab:slowfast} decomposes each sample into the two populations. The
columns headed ``in band'' are the check that the decomposition means what it
claims: the slow population's implied driving frequencies should fall in the
microseismic band and the fast population's in the anthropogenic band. If they
do not, the split is a curve-fitting artefact rather than a physical
separation.

\begin{table}[htbp]
\centering
\caption{Slow/fast decomposition. ``In band'' is the fraction of each
population whose implied driving frequency falls in the corresponding
ground-motion band.}
\label{tab:slowfast}
\small
\generated{tab_slowfast}
\end{table}

\subsection{Auxiliary correlation in O3}
\label{sec:results:aux}

\subsubsection{Rate versus ground motion}

\Cref{tab:ratecorr} gives the Spearman correlation between the hourly
scattered-light rate and the hourly band-limited RMS of each ground-motion
channel, with six-hour moving-block bootstrap intervals.

\begin{table}[htbp]
\centering
\caption{Rank correlation of the hourly scattered-light rate with
ground-motion band-limited RMS, \Lone{} \Otzb. Channel names are abbreviated:
each is prefixed \code{L1:ISI-GND\_STS\_} and suffixed \code{.rms}. Intervals
are \SI{95}{\percent} moving-block bootstrap. Full table in
\code{results/tables/aux\_rate\_correlations.csv}.}
\label{tab:ratecorr}
\footnotesize
\generated{tab_ratecorr}
\end{table}

The strongest channel is \ResultStrongestChannel{}, at
$\rho = \ResultStrongestRho$.

Three supporting checks accompany it. The top-to-bottom quartile rate ratio
for that channel is \ResultQuartileRatio{}, which states the effect in
operational terms: how much more often scattering occurs in the noisiest
quarter of hours than in the quietest. The lag scan peaks at
\ResultLagPeak{}~hours, against the zero expected for an instantaneous
mechanical response. The partial correlation of rate with the anthropogenic
band, controlling for the microseism, tests whether the two bands carry
independent information or merely track one another.

\begin{figure}[htbp]
\centering
\begin{subfigure}{0.49\textwidth}
  \safefig[\textwidth]{aux_correlation_L1_O3b}
  \caption{Correlation summary.}
\end{subfigure}\hfill
\begin{subfigure}{0.49\textwidth}
  \safefig[\textwidth]{rate_vs_groundmotion_L1_O3b}
  \caption{Rate against ground motion.}
\end{subfigure}
\caption{Left: Spearman $\rho$ with moving-block bootstrap intervals for every
ground-motion channel. Right: hourly rate against the band-limited RMS of the
strongest channel, with binned medians and interquartile range.}
\label{fig:auxcorr}
\end{figure}

\subsubsection{The per-trigger predictor}

\Cref{tab:witness} gives the fitted witness calibration and the agreement
measured on the disjoint validation set. The calibration is quoted with the
range over which the training objective stays within \SI{1}{\percent} of its
maximum: a narrow range indicates that the witness genuinely constrains the
prediction, a wide one that it does not.

\begin{table}[htbp]
\centering
\caption{Scattering predictor $f_n = 2n|v|/\lambda$ against the witness
\code{L1:SUS-OM3\_M1\_DAMP\_L\_IN1\_DQ}. The calibration $\alpha$ is fitted on
a training set and frozen before the validation measurement; the
\SI{1}{\percent} range is the interval over which the training objective stays
within one per cent of its maximum.}
\label{tab:witness}
\small
\generated{tab_witness}
\end{table}

\begin{figure}[htbp]
\centering
\safefig[0.9\textwidth]{predictor_overlay_1_L1}
\caption{A real scattered-light trigger with the harmonic comb predicted from
the auxiliary witness overlaid (dashed). The prediction uses only auxiliary
data and the single fitted scale factor; the arches are from strain. Agreement
between the two is the evidence that \cref{eq:fn} describes the mechanism.}
\label{fig:predictor}
\end{figure}

\subsection{The O4a search}
\label{sec:results:o4a}

\Cref{tab:search} records the coverage of the livetime-limited search and the
number of candidates classified as scattered light. The identifier's
cross-validated performance on held-out O3 folds is
$\mathrm{AUC} = \ResultClassifierAuc$ with an efficiency of
\ResultClassifierEfficiency{} at a false-positive rate of $0.01$ against the
control classes --- with the two caveats of \cref{sec:methods:classifier}
attached.

\begin{table}[htbp]
\centering
\caption{Blind-search coverage and yield. Rates are per hour of \emph{analysed}
livetime; the coverage fraction is the analysed livetime as a fraction of the
epoch's total.}
\label{tab:search}
\small
\generated{tab_search}
\end{table}

\subsection{O3 versus O4a}
\label{sec:results:comparison}

\subsubsection{Rate}

With the identical search applied to both runs, the ratio of the \Ofa{} rate
to the O3 rate is \ResultRateRatioHOne{} at \Hone{} and \ResultRateRatioLOne{}
at \Lone, with Poisson intervals as given in \cref{tab:search}.

\subsubsection{Composition}

The headline structural question is whether the slow, microseism-driven
component was suppressed relative to the fast component. At \Hone{} the slow
fraction went from \ResultSlowFractionOThreeHOne{} in O3 to
\ResultSlowFractionOFouraHOne{} in \Ofa, a change of
\ResultSlowShiftHOne{}; at \Lone{} it went from
\ResultSlowFractionOThreeLOne{} to \ResultSlowFractionOFouraLOne{}, a change of
\ResultSlowShiftLOne{} (\SI{95}{\percent} Newcombe intervals).

\subsubsection{Morphology}

\Cref{tab:featurecomp} lists the features whose distributions differ most
between the runs, ranked by Cliff's $\delta$ rather than by $p$-value, for the
reason given in \cref{sec:methods:stats}. The largest effect is in
\ResultStrongestFeature{}, at $\delta = \ResultStrongestDelta$.

\begin{table}[htbp]
\centering
\caption{Feature-by-feature comparison, ranked by effect size. $p$-values are
Benjamini--Hochberg corrected across all twenty features. Full table in
\code{results/tables/feature\_comparison.csv}.}
\label{tab:featurecomp}
\small
\generated{tab_featurecomp}
\end{table}

\begin{figure}[htbp]
\centering
\safefig[\textwidth]{feature_distributions_o3_vs_o4a}
\caption{Morphology feature distributions for the O3 and O4a samples,
normalised to unit area. Both samples are produced by the same search and the
same classifier, so differences reflect the detector rather than the pipeline.}
\label{fig:featuredist}
\end{figure}

\subsection{Systematics}
\label{sec:results:systematics}

Four checks bound the result.

\paragraph{Label purity.} The Gravity Spy confidence cut is varied from
$0.80$ to $0.95$. The O3 morphology medians move by \result{shift}, which
bounds the effect of label contamination.

\paragraph{Trigger generator.} The two comparison modes of
\cref{sec:methods:search} --- like-for-like against the same search on O3, and
against the Gravity Spy sample --- are reported side by side. Agreement
indicates the result does not depend on how triggers were found; disagreement
is quantified and attributed.

\paragraph{Sampling.} The livetime-limited search covers only
\ResultCoverageOFoura{} of O4a, so it is blind to structure on timescales longer than the
sampling interval. The windows are drawn uniformly in livetime and stratified
across the run, which makes the sample representative of the run as a whole
but does not recover, for example, a two-week period of unusual behaviour.

\paragraph{Extraction bias.} The \SI{15}{\percent} high bias in maximum arch
frequency (\cref{sec:validation}) applies identically to both runs and so
cancels in the comparison. It is not corrected for, and no absolute statement
is made using that feature.

\section{Discussion}
\label{sec:discussion}

\subsection{What the arch period buys}

The single most useful property of scattered light for an analysis restricted
to public data is that \cref{eq:fdrive} makes the driving frequency a
\emph{strain-only} measurement. No auxiliary channel is required to say that a
particular trigger was caused by motion at \SI{0.15}{\hertz} rather than at
\SI{2}{\hertz}; the arch period says so directly.

That is what allows the O4a half of this study to exist at all. The O4a
auxiliary release contains nothing that witnesses motion
(\cref{sec:data:aux:o4a}), so an approach that depended on auxiliary data
would simply stop at the end of O3. Reading the mechanism off the morphology
does not.

The inference is only as good as the relation it rests on, which is why
\cref{sec:results:aux} matters beyond its own numbers. Establishing in O3 ---
where both the labels and the witnesses exist --- that the arch period does
track ground motion, and that the predictor $f_n = 2n|v|/\lambda$ quantitatively
explains individual arch trains, is what licenses reading an O4a arch period
as a statement about the instrument's environment. Without the O3 auxiliary
work, the O4a arch-period distribution would be a description of pictures
rather than a measurement of physics.

\subsection{Relation to previous work}

\citet{Soni2021Scattering} established the two-population picture in O3 using
internal data, and identified relative motion between the main and reaction
chains as the dominant slow-scattering path. \citet{Nandi2026} report on
scattered light at Livingston during O4, again finding two populations: one
traced to microseismic ground motion and mitigated by installing baffles, the
other to vertical ground motion coupling through a vacuum chamber and
addressed by adding seismic isolation.

This work is not a substitute for either. It has access to a fraction of the
channels, none of the engineering record, and no ability to intervene in the
instrument. What it adds is a different kind of statement: a reproducible,
externally checkable measurement of how the scattered-light population changed
between the two runs, derived entirely from data anyone can download, with
every intermediate step validated against signals of known parameters. Where
its conclusions agree with the internal analyses, that agreement is
independent evidence; where they disagree, the disagreement is informative
about what public data can and cannot support.

The comparison with \citet{Nandi2026} is the most direct test available. If
the arch-period distribution measured here from O4a strain alone reproduces
the two-population structure they identify from internal channels, the
strain-only method is validated on real data as well as on injections.

\subsection{Limitations}

\paragraph{The O4a class definition is inherited from O3.} Every O4a number
rests on a classifier trained on O3 labels, so it presumes that the
morphological boundary of ``scattered light'' did not move between runs. That
presumption is not testable without O4a labels. The cross-validated efficiency
and false-positive rate on held-out O3 folds
(\cref{sec:results:o4a}) are the quantitative statement of how well the
transfer is expected to hold \emph{within} O3; they are an upper bound on how
well it holds across runs, not an estimate of it. A genuinely novel O4a glitch
morphology resembling scattered light would be absorbed into the sample
without comment.

\paragraph{The trigger generator is not Omicron.} The excess-power finder of
\cref{sec:methods:search} is a stand-in. Its efficiency curve differs from
Omicron's, in ways not characterised here. This is why the comparison is run
like-for-like, with the same finder on both runs: the absolute rates are
specific to this pipeline, and only the ratio between runs is meant to be
compared with anything external.

\paragraph{Coverage.} The O4a search covers a few per cent of the run. The
sample is representative in the statistical sense --- windows are drawn
uniformly in livetime and stratified across the run --- but it cannot resolve
episodic behaviour, and a fortnight of anomalous scattering would appear only
as a handful of windows.

\paragraph{One free parameter in the predictor.} The counts-to-metres
calibration is fitted, not derived, because GWOSC publishes auxiliary channels
uncalibrated. The fit is honest --- trained and validated on disjoint sets,
with the constraint width reported --- but a genuinely calibrated witness
would turn the agreement statistic from a goodness-of-fit into a real test of
\cref{eq:fn}.

\paragraph{Cross-validation is not grouped in time.} Triggers minutes apart
are not independent, so the quoted AUC is mildly optimistic. Time-grouped
folds would be the correct treatment and are the obvious next refinement.

\paragraph{\Hone{} has no public motion witness.} The per-trigger predictor is
a Livingston-only result. Whether the same calibration and the same explained
fraction would be found at Hanford is untested, and untestable from public
data.

\subsection{What would strengthen this}

In rough order of value per unit effort:

\begin{enumerate}
  \item \textbf{An O4a trigger catalogue.} Even an unlabelled Omicron trigger
        list would remove the largest systematic here, by decoupling the
        trigger generator from this pipeline.
  \item \textbf{Time-grouped cross-validation}, replacing the stratified folds
        with folds split on contiguous blocks of the run.
  \item \textbf{Full O4a coverage}, which is a matter of compute rather than
        method; the pipeline is resumable and would scale to it unchanged.
  \item \textbf{Direct comparison against \citet{Nandi2026}} on overlapping
        times, which would test the strain-only classification against an
        independent, internally-derived identification.
  \item \textbf{An O4a auxiliary release containing a motion witness}, which
        would make the O4a half of the study symmetric with the O3 half. The
        test suite is written to detect this if it happens.
\end{enumerate}

\section{Conclusions}
\label{sec:conclusions}

This work characterises scattered-light transients in LIGO O3 and O4a using
public data only, and asks how much of the physics survives that restriction.

\begin{enumerate}

\item \textbf{The morphology carries the mechanism.} Because the scattering
frequency follows $|v|$ rather than $v$, the arch repetition period measures
the driving motion directly, $\fdrive = 1/2\Tarch$, and the harmonic spacing
at the apex measures the peak surface velocity. Both are strain-only
measurements. Validated against injections, the arch period is recovered to
better than \SI{1}{\percent} and the harmonic spacing to within
\SI{7}{\percent}.

\item \textbf{Two public-data facts constrain what is possible, and neither is
documented elsewhere.} No Gravity Spy classification catalogue has been
released for O4a, and the O4a auxiliary release contains no seismometer,
accelerometer or suspension witness. Together they mean the O4a analysis must
be strain-only, and that the auxiliary correlation is necessarily an O3
result. Both were established by direct verification before any code was
written; the second is asserted in the test suite so that a future release
would be detected.

\item \textbf{The O3 auxiliary release supports more than expected.} Alongside
the complete ground-seismometer band-limited RMS family across the earthquake,
microseismic and anthropogenic bands, it contains
\code{L1:SUS-OM3\_M1\_DAMP\_L\_IN1\_DQ} at \SI{256}{\hertz} --- a genuine
optic-motion witness, and the only full-rate one in either release. This makes
a per-trigger test of $f_n = 2n|v|/\lambda$ possible from public data, which
was not obviously the case at the outset.

\item \textbf{The O3 auxiliary work is what licenses the O4a interpretation.}
Establishing in O3 that the trigger rate tracks ground motion, and that the
predictor explains individual arch trains, is what allows an O4a arch period
--- measured with no auxiliary data at all --- to be read as a statement about
the instrument's environment rather than as a description of a picture.

\item \textbf{Validation against known signals changed the analysis, twice.}
Capping the $Q$ range corrected a factor-of-twenty error in the measured arch
period of fast scattering, and raising the ridge threshold from $Z = 5.5$ to
$Z = 20$ removed a noise floor that was saturating the frequency-extent
features. Neither error was visible in the outputs; both were visible against
injections. The residual \SI{15}{\percent} high bias in maximum arch frequency
is reported rather than tuned away, and is used only where it cancels.

\end{enumerate}

The pipeline is public, resumable, and validated end to end. Its main
extension --- full O4a coverage --- is a matter of compute rather than method,
and the arrival of an O4a trigger catalogue would remove the largest remaining
systematic without changing anything else.

\subsection*{Acknowledgements}

This research has made use of data obtained from the Gravitational Wave Open
Science Center (\url{https://gwosc.org}), a service of the LIGO Laboratory, the
LIGO Scientific Collaboration, the Virgo Collaboration and KAGRA. LIGO
Laboratory and Advanced LIGO are funded by the United States National Science
Foundation, with contributions from the Science and Technology Facilities
Council of the United Kingdom, the Max-Planck-Society, and the State of
Niedersachsen/Germany.

This work uses the Gravity Spy machine-learning classification data release
\citep{GravitySpyData2021} and the GWOSC O3 and O4 auxiliary channel releases
\citep{GWOSCO3Aux,GWOSCO4Aux}. The Gravity Spy project is supported by the
United States National Science Foundation and relies on the contributions of
its citizen-science volunteers.

The analysis uses \code{gwpy} \citep{Macleod2021GWpy}, \code{numpy}
\citep{Harris2020NumPy}, \code{scipy} \citep{Virtanen2020SciPy},
\code{scikit-learn} \citep{Pedregosa2011} and \code{matplotlib}
\citep{Hunter2007}.

\subsection*{Data and code availability}

All input data are public: strain and auxiliary channels from GWOSC, and the
Gravity Spy classifications from Zenodo record 5649212. The analysis code,
configuration, verbatim channel lists and decision log are at
\url{https://github.com/Hardik-SSingh/ligo-scattered-light-o3-o4a}. Every derived
product is written with a provenance sidecar recording the configuration, the
random seed and the software versions used to produce it.


\clearpage
\bibliography{references}

\appendix
\section{The morphology feature vector}
\label{app:features}

Each Q-scan is reduced to the twenty features below, computed by the function
\code{extract\_features} in \code{gwglitch.morphology}. All are derived from the
normalised-energy map of \cref{eq:normenergy}, so none depends on the absolute
calibration and all are directly comparable between detectors and between
runs. A feature that is undefined for a given map is returned as \code{NaN}
and never as a substituted value.

Notation: $Z(t,f)$ is the normalised energy, $Z_{\mathrm{thr}} = 20$ the ridge
threshold, $f_{\mathrm{up}}(t)$ the highest supra-threshold frequency in time
bin $t$ (the arch envelope), and $E(f) = \sum_t Z(t,f)$ the supra-threshold
energy marginal.

\begin{longtable}{%
  >{\raggedright\arraybackslash}p{4.4cm}%
  >{\raggedright\arraybackslash}p{1.8cm}%
  >{\raggedright\arraybackslash}p{7.5cm}}
\toprule
Feature & Units & Definition and purpose \\
\midrule
\endfirsthead
\toprule
Feature & Units & Definition and purpose \\
\midrule
\endhead
\bottomrule
\endfoot

\multicolumn{3}{l}{\textit{Extent}}\\[0.2em]
\code{snr\_qscan} & --- &
$\sqrt{2 Z_{\max}}$, the single-tile SNR of the loudest tile. \\[0.3em]

\code{energy\_total} & --- &
Sum of supra-threshold $Z$. Spans orders of magnitude, so it is compared on a
log scale. \\[0.3em]

\code{f\_peak} & \si{\hertz} &
Frequency of the loudest tile. \\[0.3em]

\code{f\_max} & \si{\hertz} &
\SI{99}{\percent} quantile of $f_{\mathrm{up}}$. Biased high by
$\sim\SI{15}{\percent}$ because a $Q$ tile of bandwidth $f/Q$ extends above the
true apex (\cref{sec:validation}); used only where that bias cancels. \\[0.3em]

\code{f\_min} & \si{\hertz} &
\SI{1}{\percent} quantile of the lowest supra-threshold frequency. \\[0.3em]

\code{bandwidth} & \si{\hertz} &
\code{f\_max} $-$ \code{f\_min}. \\[0.3em]

\code{duration} & \si{\second} &
Total time spanned by supra-threshold bins. Separates the seconds-long
scattering trains from the millisecond blips. \\[0.5em]

\multicolumn{3}{l}{\textit{Arch structure}}\\[0.2em]
\code{n\_arches} & --- &
Number of prominent maxima of $f_{\mathrm{up}}$. \\[0.3em]

\code{arch\_period} & \si{\second} &
$\Tarch$, from the first prominent maximum of the autocorrelation of
$f_{\mathrm{up}}$ (\cref{sec:methods:features}). The physically central
feature: $\fdrive = 1/2\Tarch$. \\[0.3em]

\code{arch\_period\_scatter} & --- &
Normalised median absolute deviation of the peak-to-peak spacings. Small for a
coherent train, large for irregular structure. \\[0.5em]

\multicolumn{3}{l}{\textit{Harmonic structure}}\\[0.2em]
\code{n\_harmonics} & --- &
Number of peaks in a frequency slice through the arch apex. Measures how many
bounce orders are above threshold, and hence how non-linear the coupling is.
\\[0.3em]

\code{harmonic\_spacing} & \si{\hertz} &
Median spacing of those peaks, which equals the fundamental apex
$2 v_{\mathrm{peak}}/\lambda$ and so measures the peak surface velocity
directly. \\[0.5em]

\multicolumn{3}{l}{\textit{Shape}}\\[0.2em]
\code{ridge\_slope\_rms} & \si{\hertz\per\second} &
RMS of $\mathrm{d}f_{\mathrm{up}}/\mathrm{d}t$. \\[0.3em]

\code{ridge\_curvature\_rms} & \si{\hertz\per\second\squared} &
RMS of the second derivative. Arches are smoothly curved; whistles and lines
are not. \\[0.3em]

\code{upward\_downward\_ratio} & --- &
Ratio of rising to falling time bins along the envelope. Close to unity for a
symmetric arch train; asymmetric for classes with a fast rise and slow decay.
\\[0.5em]

\multicolumn{3}{l}{\textit{Energy distribution}}\\[0.2em]
\code{time\_symmetry} & --- &
Skewness of the energy time-marginal. \\[0.3em]

\code{frac\_energy\_below\_40hz} & --- &
$\sum_{f<40} E / \sum E$. High for slow scattering. \\[0.3em]

\code{frac\_energy\_above\_100hz} & --- &
$\sum_{f>100} E / \sum E$. Low for scattering, high for blips; one of the
strongest discriminators against the control classes. \\[0.3em]

\code{spectral\_centroid} & \si{\hertz} &
$\sum f E(f) / \sum E(f)$. \\[0.3em]

\code{band\_concentration} & --- &
$1 - W_{1/2}/W_{\mathrm{tot}}$, where $W_{1/2}$ is the width of the narrowest
frequency band containing half the energy. Near $1$ for narrow-band structure,
near $0$ for broadband. \\
\end{longtable}

\subsection*{Missingness is informative}

Several features are genuinely undefined for some maps:
\code{arch\_period} when fewer than two arches are present,
\code{harmonic\_spacing} when fewer than two harmonic peaks are resolved, and
\code{upward\_downward\_ratio} when the envelope never falls.
That a feature could not be measured is itself
diagnostic --- a blip has no arch period because it has no arches --- so the
classifier imputes the median \emph{and} carries a missingness indicator,
rather than dropping the row or silently filling a number that would be
indistinguishable from a measurement.

\section{Reproducibility}
\label{app:reproducibility}

\subsection{Environment}

The analysis runs on Python 3.11 with \code{gwpy}, \code{numpy}, \code{scipy},
\code{pandas}, \code{scikit-learn}, \code{matplotlib} and
\code{python-nds2-client}. The full specification is
\code{environment.yml}:

\begin{lstlisting}[language=bash]
conda env create -f environment.yml
conda activate gwglitch
pip install -e . --no-deps
\end{lstlisting}

Two aspects of that file are load-bearing and are recorded here because they
cost real time to diagnose.

\paragraph{OpenBLAS is pinned.} The unconstrained conda-forge solve selects
MKL 2026.1.0, in which every LAPACK call terminates the interpreter with
\code{STATUS\_STACK\_BUFFER\_OVERRUN} (exit code \code{0xC0000409}) and no
Python traceback whatsoever. \code{numpy.linalg.lstsq},
\code{numpy.linalg.solve}, \code{scipy.signal.detrend} and every
\code{scikit-learn} regression die silently. \code{environment.yml} pins the
OpenBLAS builds of \code{libblas}, \code{libcblas} and \code{liblapack}, and
the environment check exercises the backend directly so that a recurrence is
caught in a second rather than hours into a run.

\paragraph{\code{gwdetchar} is absent by design.} Its dependency chain
(\code{lalsuite}, \code{dqsegdb2}, \code{gwdatafind}) has no \code{win-64}
build on conda-forge, so it cannot be installed in a native Windows
environment. Everything the analysis needs from it --- the Q-transform and the
omega-scan rendering --- is built on \code{gwpy}'s \code{q\_transform}, which
is pure Python and SciPy. The transform is the same algorithm
\citep{Chatterji2004}; only the plotting wrapper differs.

\subsection{Verification before running}

\begin{lstlisting}[language=bash]
python scripts/00_check_environment.py
\end{lstlisting}

This checks package versions, exercises LAPACK and the SciPy filter design
directly, parses the configuration and its channel lists, and confirms that
GWOSC open strain, the NDS auxiliary server and the segment API all answer for
both O3 and O4a. It exits non-zero on any mandatory failure. The run reported
in \cref{sec:data} passed all six sections.

\subsection{Pipeline stages}

Each stage is a standalone program, resumable, and writes a JSON provenance
sidecar next to its output recording the configuration path, the random seed,
the software versions and the stage-specific parameters.

\begin{longtable}{%
  >{\raggedright\arraybackslash}p{5.7cm}%
  >{\raggedright\arraybackslash}p{8.3cm}}
\toprule
Stage & What it does \\
\midrule
\endfirsthead
\toprule
Stage & What it does \\
\midrule
\endhead
\bottomrule
\endfoot
\code{00\_check\_environment.py} &
Verifies packages, LAPACK, configuration and every remote service. \\[0.3em]

\code{01\_fetch\_gravityspy\_catalog.py} &
Downloads the four O3 CSVs ($\sim$\SI{560}{\mega\byte}) from Zenodo 5649212,
applies the selection cuts, and draws the time-stratified target and control
sub-samples. \\[0.3em]

\code{02\_run\_qscans.py} &
Q-scans each trigger and extracts the feature vector. The long stage:
network-bound at roughly one trigger per second. Checkpoints every $100$
triggers and skips completed work, so it can be interrupted freely. \\[0.3em]

\code{03\_aux\_correlation.py} &
The O3 auxiliary analysis: hourly rate against ground-motion band-limited RMS
with block-bootstrap intervals, quantile table, lag scan and partial
correlation; plus the per-trigger predictor fitted on the \Lone{} OM3 witness.
Refuses to run on an O4 epoch. \\[0.3em]

\code{04\_train\_identifier.py} &
Trains and cross-validates the strain-only classifier and fixes its operating
point. \\[0.3em]

\code{05\_search\_o4a.py} &
The livetime-limited blind search. Run once with \code{--epoch O4a} and once
with \code{--epoch O3b} for the like-for-like comparison. \\[0.3em]

\code{06\_population\_comparison.py} &
Rates, slow/fast decomposition, feature comparison; writes
\code{paper/numbers.tex}. \\[0.3em]

\code{07\_make\_figures.py} &
Renders every figure in this document to \code{results/figures}. \\
\end{longtable}

\subsection{Determinism}

Every random choice --- the stratified catalogue sub-sampling, the search
window placement, the cross-validation folds, the bootstrap resamples ---
draws from a seed set in \code{config/analysis.yaml} and recorded in each
provenance sidecar. Re-running a stage with the same configuration reproduces
its output exactly. Changing the seed changes the sample, and the size of that
change is one of the systematics of \cref{sec:results:systematics}.

\subsection{Cost}

Approximate wall-clock times on a laptop with a domestic connection:

\begin{center}
\small
\begin{tabular}{lll}
\toprule
Stage & Bottleneck & Approximate time \\
\midrule
01 catalogue download, 560 MB & network & \SIrange{5}{15}{\minute} \\
02 Q-scans, 400 blocks, $\sim$5600 events & network & \SIrange{15}{40}{\hour} \\
03 auxiliary correlation & network & \SIrange{2}{6}{\hour} \\
04 classifier training & CPU & \SIrange{2}{5}{\minute} \\
05 O4a search, 400 windows & network and CPU & \SIrange{4}{10}{\hour} \\
06 comparison & CPU & \SIrange{1}{3}{\minute} \\
07 figures & mixed & \SIrange{5}{20}{\minute} \\
\bottomrule
\end{tabular}
\end{center}

These are measured, not estimated. The dominant cost is stage~02, and it is
set by the number of GWOSC \emph{file} downloads rather than by the number of
events: open strain is served in \SI{4096}{\second} HDF5 files, and a request
for any sub-interval fetches the whole file, timed at two to four minutes each
on a domestic connection. The block cache described in
\cref{sec:methods:sampling} turns that from one download per event into one
per block; the measured yield over a real run is about ten events per
download, against a naive 1.4.

Every remote fetch is cached on disk keyed by a hash of the request, so a
re-run after an interruption costs only the work not yet done -- which matters,
because a job of this length will be interrupted. The cache is
the dominant disk consumer, at roughly \SI{1}{\mega\byte} per
\SI{64}{\second} strain window.

\subsection{Tests}

\begin{lstlisting}[language=bash]
pytest
\end{lstlisting}

$114$ tests, no network access required, about three minutes. They cover the
feature contract, the recovery tolerances of \cref{tab:validation}, the
pure-noise null test, the witness-calibration round trip, a negative control
on the predictor, the GPS conversions against the published epoch boundaries,
the segment arithmetic, the statistical estimators, and the two
data-availability facts of \cref{sec:data} --- so that a future GWOSC release
adding an O4a motion witness would announce itself as a test failure.

\subsection{Building this document}

The \code{paper} directory compiles standalone on Overleaf with pdf\LaTeX{}
and Bib\TeX{}. \code{\textbackslash graphicspath} declares both
\code{../results/figures} and \code{./figures}, so the source works whether it
sits inside the analysis repository or has been uploaded on its own; in the
latter case, copy \code{results/figures/*.pdf} into \code{paper/figures}.

Numbers awaiting the production run appear as \code{\textbackslash result\{\}}
placeholders, rendered in red so that no unfilled value can be mistaken for a
measurement. Stage~06 closes the loop in two ways. It writes
\code{paper/numbers.tex}, which redefines the inline macros to the measured
values and removes the status notice at the head of \cref{sec:results}; and it
writes \code{paper/generated/*.tex}, the body rows of each results table,
which the section files pull in through \code{\textbackslash generated} when
they exist and fall back to placeholders when they do not. The column
specifications, captions and prose stay under human control; only the numbers
are machine-written, and they cannot drift from the analysis that produced
them.

The macro names are generated mechanically from the pipeline's keys, so
\code{slow\_fraction\_o4a\_h1} becomes
\code{\textbackslash ResultSlowFractionOFouraHOne} --- digits are spelled out
because \TeX{} forbids them in a control sequence.
\code{tests/test\_latex.py} asserts that every key the pipeline emits has a
declaration in \code{main.tex} and that every macro used in a section is
declared, so the two halves cannot silently fall out of step.


\end{document}
